\section{Methodology}
\label{sec:methodology}

\subsection{Problem Formulation}
\label{sec:formulation}

We formulate the MSSP as a multi-objective mixed-integer program over a hierarchical location structure, capturing the full complexity of real-world film production logistics.

\paragraph{Primary Sets.}
Let $\mathcal{S} = \{s_1, \ldots, s_n\}$ be the set of scenes, $\mathcal{A} = \{a_1, \ldots, a_m\}$ the actors, $\mathcal{D} = \{1, \ldots, T\}$ the planning horizon of available shooting days, and $\mathcal{P} \subseteq \mathcal{S} \times \mathcal{S}$ the set of precedence pairs.

\paragraph{Hierarchical Location Model.}
Real film productions operate across a three-tier geographic hierarchy: \emph{cities} (production bases), \emph{towns} (districts within cities), and \emph{locations} (specific shooting venues within towns).
We model this via three sets and two surjective mappings:
\begin{itemize}
    \item $\mathcal{C} = \{c_1, \ldots, c_p\}$: cities (production hubs);
    \item $\mathcal{I} = \{i_1, \ldots, i_q\}$: towns, with $\mathcal{I}_c \subseteq \mathcal{I}$ denoting the towns in city $c$;
    \item $\mathcal{L} = \{l_1, \ldots, l_k\}$: locations, with $\mathcal{L}_i \subseteq \mathcal{L}$ denoting the locations in town $i$;
    \item $\phi: \mathcal{L} \to \mathcal{I}$: maps each location to its town, $\phi(l) = i \iff l \in \mathcal{L}_i$;
    \item $\psi: \mathcal{I} \to \mathcal{C}$: maps each town to its city, $\psi(i) = c \iff i \in \mathcal{I}_c$.
\end{itemize}
This induces nested containment: $l \in \mathcal{L}_i \Rightarrow i \in \mathcal{I}_{\psi(i)}$, ensuring that activating a location implies its town and city are also active.

\paragraph{Scene Parameters.}
Each scene $s \in \mathcal{S}$ has:
shooting duration $\delta_s$ (hours);
required actors $A(s) \subseteq \mathcal{A}$;
compatible locations $L(s) \subseteq \mathcal{L}$;
optional time window $[e_s, l_s] \subseteq \mathcal{D}$;
optional soft deadline $\bar{d}_s$ with per-day penalty $\rho_s$.

\paragraph{Actor Parameters.}
Each actor $a \in \mathcal{A}$ has:
daily wage $w_a$;
availability calendar $\mathcal{D}_a \subseteq \mathcal{D}$;
criticality weight $\kappa^A_a \in \mathbb{R}_{>0}$ (reflecting contractual priority);
maximum consecutive working days $\bar{h}_a$ (labor regulations).

\paragraph{Location and Transfer Parameters.}
Each location $l \in \mathcal{L}$ has:
availability calendar $\mathcal{D}_l \subseteq \mathcal{D}$;
criticality weight $\kappa^L_l \in \mathbb{R}_{>0}$ (reflecting booking difficulty).
Transfer costs capture three tiers:
intra-town $\tau^{\text{loc}}_{l,l'}$,
inter-town $\tau^{\text{town}}_{i,i'}$,
inter-city $\tau^{\text{city}}_{c,c'}$.
The composite transfer cost is:
\begin{equation}
c_{l,l'} = \begin{cases}
    \tau^{\text{loc}}_{l,l'} & \text{if } \phi(l) = \phi(l') \\[2pt]
    \tau^{\text{town}}_{\phi(l),\phi(l')} + \tau^{\text{loc}}_{l,l'} & \text{if } \psi(\phi(l)) = \psi(\phi(l')),\; \phi(l) \neq \phi(l') \\[2pt]
    \tau^{\text{city}}_{\psi(\phi(l)),\psi(\phi(l'))} + \tau^{\text{town}}_{\phi(l),\phi(l')} + \tau^{\text{loc}}_{l,l'} & \text{otherwise}
\end{cases}
\label{eq:transfer_cost}
\end{equation}

\paragraph{Binary Indicator Parameters.}
For notational convenience we define:
$b_{a,s} = \mathbb{1}[a \in A(s)]$ (actor $a$ required for scene $s$);
$\alpha_{a,d} = \mathbb{1}[d \in \mathcal{D}_a]$ (actor $a$ available on day $d$);
$\lambda_{l,d} = \mathbb{1}[d \in \mathcal{D}_l]$ (location $l$ available on day $d$);
$p_{s,l} = \mathbb{1}[l \in L(s)]$ (scene $s$ compatible with location $l$);
$\eta_{i,l} = \mathbb{1}[l \in \mathcal{L}_i]$ (location $l$ belongs to town $i$).

\paragraph{Decision Variables.}
The model uses five families of binary decision variables:
\begin{align}
x_{s,d} &\in \{0,1\} && \text{scene } s \text{ scheduled on day } d \label{var:x}\\
y_{s,l} &\in \{0,1\} && \text{scene } s \text{ assigned to location } l \label{var:y}\\
v_{a,s,d} &\in \{0,1\} && \text{actor } a \text{ participates in scene } s \text{ on day } d \label{var:v}\\
z_{l,s,d} &\in \{0,1\} && \text{scene } s \text{ shot at location } l \text{ on day } d \label{var:z}\\
u_{i,d} &\in \{0,1\} && \text{town } i \text{ is the active shooting town on day } d \label{var:u}
\end{align}
The linking variables $v_{a,s,d}$ and $z_{l,s,d}$ couple scene scheduling with actor participation and location usage, enabling fine-grained constraint enforcement.
The town activation variable $u_{i,d}$ enforces geographic crew logistics.

\paragraph{Bi-Objective Formulation.}
We minimize two conflicting objectives simultaneously:
\begin{align}
f_1(\mathbf{x}, \mathbf{v}, \mathbf{z}) &= \underbrace{\sum_{a \in \mathcal{A}} w_a \cdot \bigl|\{d \in \mathcal{D} : \textstyle\sum_{s} v_{a,s,d} \geq 1\}\bigr|}_{\text{(i) actor daily wages}} \;+\; \underbrace{\sum_{d=2}^{T} \sum_{\substack{l, l' \in \mathcal{L} \\ l \neq l'}} c_{l,l'} \cdot \sigma_{l,l',d}}_{\text{(ii) hierarchical transfer costs}} \notag \\[4pt]
&\quad +\; \underbrace{\sum_{s \in \mathcal{S}} \rho_s \cdot \max\bigl(0,\, \textstyle\sum_{d} d \cdot x_{s,d} - \bar{d}_s\bigr)}_{\text{(iii) deadline penalties}} \;+\; \underbrace{\sum_{(l,s,d)} \frac{z_{l,s,d}}{\kappa^L_l} \;+\; \sum_{(a,s,d)} \frac{v_{a,s,d}}{\kappa^A_a}}_{\text{(iv) criticality-weighted utilization}} \label{eq:cost}\\[8pt]
f_2(\mathbf{x}) &= \max_{s \in \mathcal{S}} \textstyle\sum_d d \cdot x_{s,d} \;-\; \min_{s \in \mathcal{S}} \textstyle\sum_d d \cdot x_{s,d} + 1 \label{eq:makespan}
\end{align}
where $\sigma_{l,l',d} = 1$ if location $l'$ is used on day $d$ and location $l$ was used on day $d{-}1$.
Cost objective $f_1$ has four terms: (i)~actor wages paid per active day, (ii)~hierarchical transfer costs from Eq.~\eqref{eq:transfer_cost}, (iii)~deadline penalties, and (iv)~criticality-weighted utilization that penalizes usage of low-criticality resources more heavily, encouraging efficient allocation of high-value actors and premium locations.
Makespan $f_2$ minimizes total production duration.

\paragraph{Constraints.}
The formulation enforces 13 constraint families organized into four categories.

\medskip
\noindent\textbf{I. Assignment Constraints.}
\vspace{-4pt}
\begin{align}
\sum_{d \in \mathcal{D}} x_{s,d} &= 1 && \forall s \in \mathcal{S} \tag{C1}\label{c1}\\
\sum_{l \in L(s)} y_{s,l} &= 1 && \forall s \in \mathcal{S} \tag{C2}\label{c2}\\
\text{day}(s_i) + 1 &\leq \text{day}(s_j) && \forall (s_i, s_j) \in \mathcal{P} \tag{C3}\label{c3}\\
e_s \leq \text{day}(s) &\leq l_s && \forall s \text{ with time windows} \tag{C4}\label{c4}
\end{align}
\ref{c1}~schedules each scene exactly once.
\ref{c2}~assigns each scene to exactly one compatible location.
\ref{c3}~enforces precedence with a minimum one-day separation.
\ref{c4}~restricts scheduling within time windows, where $\text{day}(s) = \sum_{d} d \cdot x_{s,d}$.

\medskip
\noindent\textbf{II. Capacity and Availability Constraints.}
\vspace{-4pt}
\begin{align}
\sum_{s \in \mathcal{S}} \delta_s \cdot x_{s,d} &\leq W_{\max} && \forall d \in \mathcal{D} \tag{C5}\label{c5}\\
x_{s,d} &\leq \alpha_{a,d} && \forall s \in \mathcal{S},\; a \in A(s),\; d \in \mathcal{D} \tag{C6}\label{c6}\\
z_{l,s,d} &\leq \lambda_{l,d} && \forall l \in L(s),\; s \in \mathcal{S},\; d \in \mathcal{D} \tag{C7}\label{c7}
\end{align}
\ref{c5}~caps daily shooting at $W_{\max}{=}8$ hours.
\ref{c6}~prevents scheduling a scene when any required actor is unavailable.
\ref{c7}~ensures locations are used only on available days.

\medskip
\noindent\textbf{III. Variable Linking Constraints.}
\vspace{-4pt}
\begin{align}
v_{a,s,d} &= x_{s,d} \cdot b_{a,s} \cdot \alpha_{a,d} && \forall a \in \mathcal{A},\; s \in \mathcal{S},\; d \in \mathcal{D} \tag{C8}\label{c8}\\
z_{l,s,d} &= x_{s,d} \cdot p_{s,l} \cdot \lambda_{l,d} && \forall l \in \mathcal{L},\; s \in \mathcal{S},\; d \in \mathcal{D} \tag{C9}\label{c9}
\end{align}
\ref{c8}~forces $v_{a,s,d} = 1$ iff scene $s$ is on day $d$, actor $a$ is required, and $a$ is available.
\ref{c9}~forces $z_{l,s,d} = 1$ iff scene $s$ is on day $d$, location $l$ is compatible, and $l$ is available.
Since $b_{a,s}$, $\alpha_{a,d}$, $p_{s,l}$, $\lambda_{l,d}$ are fixed binary parameters, these reduce to linear equalities.

\medskip
\noindent\textbf{IV. Geographic Logistics Constraints.}
\vspace{-4pt}
\begin{align}
z_{l,s,d} &\leq u_{\phi(l),d} \cdot \eta_{\phi(l),l} \cdot \lambda_{l,d} && \forall l \in L(s),\; s \in \mathcal{S},\; d \in \mathcal{D} \tag{C10}\label{c10}\\
\sum_{i \in \mathcal{I}} u_{i,d} &\leq 1 && \forall d \in \mathcal{D} \tag{C11}\label{c11}\\
u_{i,d} + u_{j,d-1} &\leq 1 && \forall i, j \in \mathcal{I},\; i \neq j,\; d \geq 2 \tag{C12}\label{c12}\\
\sum_{d'=d}^{\min(d + \bar{h}_a - 1,\, T)} \!\!\Bigl(\sum_{s \in \mathcal{S}} v_{a,s,d'}\Bigr) &\leq \bar{h}_a && \forall a \in \mathcal{A},\; d \in \mathcal{D} \tag{C13}\label{c13}
\end{align}
\ref{c10}~enforces location--town consistency: using location $l$ on day $d$ requires its town $\phi(l)$ to be active.
\ref{c11}~restricts the crew to \emph{at most one active town per day}---a core logistics constraint, as crew and equipment cannot operate in multiple geographic areas simultaneously.
\ref{c12}~enforces \emph{town continuity across consecutive days}: the crew cannot switch towns between successive shooting days, capturing the real-world constraint that overnight location changes are infeasible without a dedicated travel day.
\ref{c13}~limits each actor to at most $\bar{h}_a$ consecutive working days, modeling labor regulations and union rules.

\paragraph{Formulation Complexity.}
The model comprises $O(nmT + nkT + qT)$ binary variables and $O(nmT + nkT + q^2 T + |\mathcal{P}|)$ constraints, where $n{=}|\mathcal{S}|$, $m{=}|\mathcal{A}|$, $k{=}|\mathcal{L}|$, $q{=}|\mathcal{I}|$, $T{=}|\mathcal{D}|$.
The linking constraints (\ref{c8}--\ref{c9}) and geographic logistics constraints (\ref{c10}--\ref{c12}) create dense variable coupling that makes direct MILP solution intractable for $n > 20$, motivating the decomposition approach in the following subsections.

\subsection{Scene Interaction Graph (SIG)}
\label{sec:sig}

We define an undirected weighted graph $G = (\mathcal{S}, E, w)$ where each scene is a node and edge weights capture structural coupling.
The SIG leverages the hierarchical location model (Section~\ref{sec:formulation}) by computing transfer costs $\bar{c}(s_i, s_j)$ using the composite three-tier formula of Eq.~\eqref{eq:transfer_cost}:
\begin{equation}
w(s_i, s_j) = \alpha \cdot J_A(s_i, s_j) + \beta \cdot J_L(s_i, s_j) + \gamma \cdot \frac{1}{1 + \bar{c}(s_i, s_j)}
\label{eq:sig}
\end{equation}
where $J_A(s_i, s_j) = \frac{|A(s_i) \cap A(s_j)|}{|A(s_i) \cup A(s_j)|}$ is the Jaccard similarity of actor sets, $J_L$ is the Jaccard similarity of location sets, $\bar{c}(s_i, s_j)$ is the average composite transfer cost between compatible locations (reflecting city/town/location distances), and $\alpha + \beta + \gamma = 1$ are weighting parameters (default: $\alpha{=}0.5$, $\beta{=}0.3$, $\gamma{=}0.2$).

Scenes sharing many actors and co-located within the same town receive high edge weights, encoding the intuition that such scenes benefit from consecutive scheduling and naturally satisfy the town-uniqueness constraint~\ref{c11}.

\paragraph{Louvain Decomposition.}
We apply Louvain community detection~\citep{blondel2008} to $G$, which maximizes modularity to identify densely connected scene clusters.
The resulting communities form \emph{shooting blocks} $B_1, B_2, \ldots, B_K$, where each block contains tightly coupled scenes.
We iteratively increase the resolution parameter and split blocks exceeding a maximum size to ensure at least $\lceil n/4 \rceil$ blocks, providing sufficient combinatorial richness for the upper-level search.

\subsection{Bi-Level Matheuristic}
\label{sec:bilevel}

\paragraph{Upper Level: NSGA-II over Block Permutations.}
The upper level encodes solutions as permutations of block indices $\pi = (\pi_1, \pi_2, \ldots, \pi_K)$, defining the order in which blocks are scheduled.
This reduces the search space from $n!$ (scene permutations) to $K!$ (block permutations), where typically $K \ll n$.

We use NSGA-II~\citep{deb2002} with population size 20 and 40 generations.
Non-dominated sorting partitions the population into Pareto fronts, and crowding distance maintains solution diversity within each front.

\paragraph{Lower Level: Global Greedy Reassembly.}
For each block permutation $\pi$, we construct a scene-level ordering by concatenating the scenes of each block in sequence: first all scenes from $B_{\pi_1}$ (sorted by constraint tightness), then $B_{\pi_2}$, and so on.
A global greedy decoder then processes this ordering, packing scenes into shooting days while enforcing all 13 constraint families from Section~\ref{sec:formulation}:
\begin{itemize}
    \item \textbf{Assignment \& Capacity} (\ref{c1}--\ref{c5}): each scene assigned to exactly one day, respecting capacity ($\leq W_{\max}$ hours/day), precedence, and time windows;
    \item \textbf{Availability} (\ref{c6}--\ref{c7}): actor and location availability calendars are checked before placement;
    \item \textbf{Geographic logistics} (\ref{c10}--\ref{c12}): location is chosen to respect town-uniqueness (at most one active town per day) and town-continuity (no town switches on consecutive days), with transfer cost minimized using Eq.~\eqref{eq:transfer_cost};
    \item \textbf{Labor limits} (\ref{c13}): actor consecutive working-day limits are verified before scheduling;
    \item Days are shared across blocks naturally, avoiding the cost overhead of isolated per-block scheduling.
\end{itemize}

This \emph{global reassembly} approach ensures that the decomposition guides the ordering while the greedy decoder handles day packing with full constraint enforcement, achieving cost-competitive solutions while preserving the search-space reduction of block-level evolution.

\paragraph{Solution Evaluation.}
Each chromosome $\pi$ is evaluated by the global greedy reassembly: scenes are ordered according to the block permutation, then greedily assigned to days.
Evaluation results are cached (keyed by $\pi$) to avoid redundant computation.
The result is a bi-objective vector $(f_1, f_2)$ fed to NSGA-II.

\begin{algorithm}[t]
\caption{BDM-AOS: Bi-Level Decomposition Matheuristic}
\label{alg:bdm_aos}
\begin{algorithmic}[1]
\REQUIRE MSSP instance $(\mathcal{S}, \mathcal{A}, \mathcal{C}, \mathcal{I}, \mathcal{L}, \mathcal{D}, \phi, \psi)$, parameters $(N, G, p_c, p_m)$
\ENSURE Pareto front of schedules
\STATE Build Scene Interaction Graph $G_{\text{SIG}}$ via Eq.~(\ref{eq:sig})
\STATE $\{B_1, \ldots, B_K\} \leftarrow$ Louvain$(G_{\text{SIG}})$
\STATE Initialize population $P = \{\pi_1, \ldots, \pi_N\}$ (random block permutations)
\STATE Evaluate each $\pi_i$ via global greedy reassembly
\FOR{$g = 1$ to $G$}
    \STATE Compute diversity $d$, improvement rate $r$, quality gap $q$
    \STATE $(cx, mut) \leftarrow$ AOS.select$(g, G, d, r, q)$ \COMMENT{Q-learning}
    \STATE Generate offspring via $cx$ crossover and $mut$ mutation
    \STATE Evaluate offspring via bi-level solving
    \STATE $P \leftarrow$ NSGA-II environmental selection$(P \cup \text{offspring})$
    \STATE AOS.update(reward $= \Delta$HV)
\ENDFOR
\RETURN Non-dominated solutions from $P$
\end{algorithmic}
\end{algorithm}

\subsection{Q-Learning Adaptive Operator Selection}
\label{sec:aos}

We formulate operator selection as a reinforcement learning problem.

\paragraph{State Space.}
The state is discretized from four continuous features:
(i)~generation progress (early/mid/late),
(ii)~population diversity via normalized Hamming distance (low/medium/high),
(iii)~improvement rate of hypervolume over the last generation (stagnant/slow/fast),
(iv)~quality gap from the reference point (good/moderate/poor).
This yields $3 \times 3 \times 3 \times 3 = 81$ possible states.

\paragraph{Action Space.}
Each action selects a (crossover, mutation) pair from $\{$PMX, OX, LAX$\} \times \{$swap, block-swap, ACM$\}$, giving 9 actions.

\paragraph{Reward.}
The reward after applying action $a$ in state $s$ is the change in hypervolume:
$r = \text{HV}(P_{g+1}) - \text{HV}(P_g)$.

\paragraph{Update Rule.}
Standard Q-learning with $\epsilon$-greedy exploration ($\epsilon_0 = 0.3$, decay $0.995$, minimum $0.05$):
\begin{equation}
Q(s, a) \leftarrow Q(s, a) + \alpha \left[ r + \gamma \max_{a'} Q(s', a') - Q(s, a) \right]
\label{eq:qlearning}
\end{equation}
with learning rate $\alpha = 0.1$ and discount factor $\gamma = 0.9$.

\subsection{Novel Problem-Specific Operators}
\label{sec:novel_ops}

\paragraph{Location-Aware Crossover (LAX).}
LAX preserves \emph{location runs}: maximal sub-sequences of blocks sharing the same primary location in a parent.
The longest location run from parent 1 is copied intact to the child, and remaining positions are filled from parent 2 in order.
This preserves spatial locality, reducing inter-block transfer costs.

\paragraph{Actor-Continuity Mutation (ACM).}
ACM identifies the actor with the largest idle gap between block appearances.
It moves the later block immediately after the earlier one, reducing the number of idle days (and thus wage payments) for that actor.
This directly exploits the MSSP cost structure where actor wages dominate total cost.

\subsection{Computational Complexity Analysis}
\label{sec:complexity}

\paragraph{Problem Complexity.}
The MSSP is NP-hard by reduction from the single-machine weighted completion time problem~\citep{garey1979, pinedo2016}.
Even the single-objective variant (minimizing cost only) requires searching over $n!$ scene orderings subject to the 13 constraint families---including the geographic logistics constraints (\ref{c10}--\ref{c12}) and actor labor limits (\ref{c13})---that create dense variable coupling with $O(nmT + nkT + qT)$ binary variables.

\paragraph{Search Space Reduction.}
BDM-AOS exploits the SIG structure to decompose $n$ scenes into $K$ blocks, reducing the upper-level search space from $O(n!)$ to $O(K!)$.
For our largest instance (S100: $n=100$ scenes), Louvain decomposition yields $K \approx 24$ blocks, reducing the search space by a factor exceeding $\frac{100!}{24!} \approx 10^{134}$.
Even at moderate scale (S50: $n=50$, $K \approx 12$), the reduction factor is $\frac{50!}{12!} \approx 10^{55}$.

\paragraph{Per-Generation Complexity.}
Each generation of the NSGA-II upper level involves:
\begin{itemize}
    \item \textit{Crossover/mutation}: $O(N \cdot K)$ per offspring, where $N$ is population size and $K$ is the number of blocks;
    \item \textit{Greedy reassembly decoding}: $O(n \cdot T)$ per individual, where $T$ is the planning horizon;
    \item \textit{Non-dominated sorting}: $O(N^2 \cdot 2)$ for bi-objective problems;
    \item \textit{AOS Q-table lookup}: $O(1)$ (81-state, 9-action table).
\end{itemize}

The total per-generation cost is $O(N \cdot (n \cdot T + K))$, dominated by the greedy decoder evaluations.

\paragraph{Overall Complexity.}
With $G$ generations and population size $N$:
$\text{Total} = O(G \cdot N \cdot n \cdot T)$.
Using our default parameters ($N=20$, $G=40$), this is $O(800 \cdot n \cdot T)$, linear in instance size for fixed parameters---a significant advantage over exact methods whose worst-case complexity is exponential.
